%% kapa: exam=1 round=21 year=2010 period=2 form=A subject=appraisal_law subject_ko=부동산관계법규 seq=36 number=76
%% source: Q-Net 감정평가사 기출문제 (cst00309) artlSeq=5024181 file=extracted/1008/5024181/2010년 제21회 감정평가사 1차시험 기출문제/2교시.hwp
%% conversion: kordoc 4.24.51 -> tools/gen1.py (auto); review=proofread-gpt-sol; flags=none
\begin{question}{76}
측량ㆍ수로조사 및 지적에 관한 법령상 토지소유자의 정리에 관한 설명으로 옳지 않은 것은?
\begin{choices}
  \item 지적공부에 신규등록하는 토지의 소유자는 등기관이 직접 조사하여 등록한다.
  \item 지적공부에 등록된 토지소유자의 변경사항은 등기관서에서 등기한 것을 증명하는 등기필 통지서, 등기필증, 등기부 등본․초본 또는 등기관서에서 제공한 등기전산정보자료에 따라 정리한다.
  \item 「국유재산법」에 따라 소유자 없는 부동산에 대한 소유자 등록을 총괄청이나 관리청이 신청하는 경우, 지적소관청은 지적공부에 해당 토지의 소유자가 등록되지 아니한 경우에만 등록할 수 있다.
  \item 등기부에 적혀 있는 토지의 표시가 지적공부와 일치하지 아니하여 토지소유자를 정리할 수 없는 경우 토지의 표시와 지적공부가 일치하지 아니하다는 사실을 관할 등기관서에 통지하여야 한다.
  \item 지적소관청은 필요하다고 인정하는 경우에는 관할 등기관서의 등기부를 열람하여 지적공부와 부동산등기부가 일치하는지 여부를 조사․확인하여야 한다.
\end{choices}
\end{question}
\answer{1}
